% Source: source/普通高考/2018/2018天津文.pdf, page 1, question 4.
% pdfimages -list found no embedded image objects; the source flowchart is vector artwork.
% Grid evidence: tmp/redraw_flowchart_2018_tianjin_liberal/recheck_source/page1_grid100.jpg;
% the original-side comparison crop is from the no-grid page render.
% Elements: rounded start/end nodes, input/output parallelograms, an initialization
% rectangle, an increment-operation diamond, two decision diamonds, directed
% connectors, a loop-back branch, and Chinese 是/否 labels. This is a schematic
% program flowchart, not a scale drawing.
% Relations: 开始 -> 输入 N -> i=2,T=0 -> test N/i integer; the 是 branch increments
% T, both branches reach i=i+1, then test i>=5; 否 loops to the first test, 是 outputs T.
% Node-edge list: start->input; input->init; init->testdiv; testdiv(是)->add;
% testdiv(否)->inc; add->inc; inc->check; check(是)->output; check(否)->testdiv;
% output->finish.  No other connectors are present.
% solid segments: every node border, vertical/horizontal connector, and arrow shaft.
% dashed segments: none.
% labels: formulas are centered in nodes; 是/否 sit beside their corresponding exits,
% clear of connector lines; return arrows enter the top of the first decision diamond.

\begin{tikzpicture}[
    >=Stealth,
    x=1cm,
    y=0.88cm,
    line width=0.45pt,
    line cap=round,
    line join=round,
    font=\normalsize,
    every node/.style={anchor=center, align=center,
        execute at begin node={\setlength{\parindent}{0pt}}},
    flowbox/.style={draw, anchor=center, minimum height=0.68cm, inner xsep=4pt, inner ysep=2pt, align=center,
        execute at begin node={\setlength{\parindent}{0pt}}},
    terminal/.style={flowbox, rounded corners=0.14cm, minimum width=1.25cm},
    io/.style={flowbox, trapezium, minimum height=0.76cm,
        minimum width=1.15cm, trapezium left angle=70, trapezium right angle=110},
    process/.style={flowbox, minimum width=2.65cm},
    decision/.style={flowbox, diamond, aspect=4.4, minimum width=3.65cm,
        minimum height=1.10cm},
]
    % Preserve white margins around both outer return paths and branch labels.
    \path[use as bounding box] (-3.10,-0.42) rectangle (3.15,11.42);

    % Coordinates preserve the source topology and relative branch layout.
    \node[terminal] (start) at (0,11.0) {开始};
    \node[io, font=\small] (input) at (0,9.72) {输入\(N\)};
    \node[process, minimum width=2.65cm] (init) at (0,8.48) {$i=2,T=0$};
    \node[decision, minimum width=4.00cm, minimum height=1.20cm] (testdiv)
        at (0,7.02) {$\displaystyle\frac{N}{i}\text{是整数?}$};
    \node[process, minimum width=1.82cm] (add) at (0,5.42) {$T=T+1$};
    \node[decision, minimum width=2.70cm, minimum height=0.90cm] (inc)
        at (0,4.05) {$i=i+1$};
    \node[decision, minimum width=2.55cm, minimum height=0.90cm] (check)
        at (0,2.63) {$i\geq5$};
    \node[io, minimum width=1.20cm, font=\small] (output) at (0,1.25) {输出\(T\)};
    \node[terminal] (finish) at (0,0.02) {结束};

    % Main downward path and the integer-test yes branch.
    \draw[->] (start.south) -- (input.north);
    \draw[->] (input.south) -- (init.north);
    \draw (init.south) -- (0,7.80);
    \draw[->] (0,7.80) -- (testdiv.north);
    \draw[->] (testdiv.south) -- node[right=2pt] {是} (add.north);
    \draw[->] (add.south) -- (inc.north);
    \draw[->] (inc.south) -- (check.north);
    \draw[->] (check.south) -- node[right=2pt] {是} (output.north);
    \draw[->] (output.south) -- (finish.north);

    % Integer-test no branch goes right and merges above i=i+1.
    \coordinate (divNoTurn) at (2.35,7.02);
    % Leave a clear vertical gap between the horizontal merge arrow and the
    % main downward arrow into the increment decision.
    \coordinate (divNoMerge) at (2.35,4.60);
    \coordinate (divNoJoin) at (0,4.60);
    \draw (testdiv.east) -- (divNoTurn) -- (divNoMerge);
    \draw[->] (divNoMerge) -- (divNoJoin);
    \draw (divNoJoin) -- (inc.north);
    \node[anchor=south west, inner sep=1pt] at (1.82,7.16) {否};

    % i>=5 no branch returns left and merges into the first decision's trunk.
    \coordinate (loopLeft) at (-2.55,2.63);
    % Separate the horizontal return arrowhead from the main-chain arrowhead.
    \coordinate (loopTop) at (-2.55,7.80);
    \coordinate (loopJoin) at (0,7.80);
    \draw (check.west) -- (loopLeft) -- (loopTop);
    \draw (loopTop) -- (loopJoin);
    \draw (loopJoin) -- (testdiv.north);
    % Place the label beside the check diamond's left exit, not on the outer spine.
    \node[anchor=south east, inner sep=1pt] at (-1.40,2.82) {否};
\end{tikzpicture}
